\chapter{Bijektionen aus gegebenen Familien}
\label{b08guide:ch:constructions}

Die folgenden Konstruktionen setzen eine \emph{gegebene Familie}
von Bijektionen voraus. Die bloße Aussage, daß für jeden Index
irgendeine Bijektion existiert, ersetzt diese Voraussetzung nicht.
Alle Aussagen gelten in ZF ohne Auswahlaxiom.

\section{Markierte Vereinigungen und Produkte}

\begin{definition}[Markierte Vereinigung und Mengenprodukt]
\label{b08guide:def:family-constructions}
Für eine Familie $(A_i)_{i\in I}$ von Mengen seien
\[
 \begin{aligned}
 \coprod_{i\in I}A_i
   &=\{(i,a)\mid i\in I,\ a\in A_i\},\\
 \prod_{i\in I}A_i
   &=\{x\colon I\to\bigcup_{i\in I}A_i\mid
                \forall i\in I\quad x(i)\in A_i\}.
 \end{aligned}
\]
Die Markierung hält verschiedene Komponenten auch bei
überlappenden Mengen auseinander. Für $I=\varnothing$ ist
die markierte Vereinigung leer; das Produkt besteht aus genau
der leeren Funktion.
\end{definition}

\begin{theorem}[Markierte Vereinigung von Bijektionen]
\label{b08guide:thm:tagged-union}
Sei $(f_i\colon A_i\bij B_i)_{i\in I}$ eine gegebene Familie
von Bijektionen. Dann ist
\[
 F\colon\coprod_{i\in I}A_i\bij\coprod_{i\in I}B_i,
 \qquad F(i,a)=(i,f_i(a)),
\]
eine Bijektion. Ihre Umkehrung lautet
\[
 F^{-1}(i,b)=(i,f_i^{-1}(b)).
\]
Leere Komponenten sind zulässig. Bei leerer Indexmenge ist
$F$ die leere Bijektion.
\end{theorem}

\begin{theorem}[Beliebiges Produkt gegebener Bijektionen]
\label{b08guide:thm:family-product}
Sei $(f_i\colon A_i\bij B_i)_{i\in I}$ eine gegebene Familie
von Bijektionen. Dann ist
\[
 P\colon\prod_{i\in I}A_i\bij\prod_{i\in I}B_i,
 \qquad P(x)(i)=f_i(x(i)),
\]
eine Bijektion mit
\[
 P^{-1}(y)(i)=f_i^{-1}(y(i))\qquad(i\in I).
\]
Eine Nichtleere der Produkte wird nicht vorausgesetzt und nicht
behauptet. Ist eines der beiden Produkte leer, so auch das andere.
Bei leerem $I$ ist $P$ die Identität auf der einelementigen Menge,
die nur die leere Funktion enthält.
\end{theorem}

\section{Bijektiver Wechsel der Indexmenge}

\begin{theorem}[Umindizierung und gleichzeitiger Komponentenwechsel]
\label{b08guide:thm:reindexing}
Seien $\sigma\colon I\bij J$ eine Bijektion, $(A_i)_{i\in I}$ und
$(B_j)_{j\in J}$ Mengenfamilien sowie
$(f_i\colon A_i\bij B_{\sigma(i)})_{i\in I}$ eine gegebene Familie
von Bijektionen. Dann sind
\[
 \begin{aligned}
 S\colon\coprod_{i\in I}A_i&\bij\coprod_{j\in J}B_j,
 & S(i,a)&=(\sigma(i),f_i(a)),\\
 T\colon\prod_{i\in I}A_i&\bij\prod_{j\in J}B_j,
 & T(x)(j)&=f_{\sigma^{-1}(j)}\bigl(x(\sigma^{-1}(j))\bigr)
 \end{aligned}
\]
Bijektionen. Ihre Umkehrungen sind durch
\[
 \begin{aligned}
 S^{-1}(j,b)
   &=\bigl(\sigma^{-1}(j),f_{\sigma^{-1}(j)}^{-1}(b)\bigr),\\
 T^{-1}(y)(i)&=f_i^{-1}\bigl(y(\sigma(i))\bigr)
 \end{aligned}
\]
gegeben. Für $B_j=A_{\sigma^{-1}(j)}$ und
$f_i=\operatorname{id}_{A_i}$ erhält man die reine Umindizierung.
Die Aussagen schließen leere Indexmengen, leere Komponenten und
leere Produkte ein.
\end{theorem}
